% 原书分章PDF第9—11页：第30.4节全部8段。
\section{物理学、统计学与科学哲学}

% S30.4.P01
一个世纪以前，物理学与哲学之间的关系仍然相对和谐，尽管相对论与量子力学等根本性革命，已经强烈动摇了传统而令人安心的世界观，使时间、空间和因果性这些古老而熟悉的哲学概念变得成问题。爱因斯坦、海森堡、玻尔（Bohr）和薛定谔（Schrödinger）等物理学家具有浓厚的哲学取向，其中许多人通过言论和著作表达了这种兴趣。罗素（Russell）、卡尔纳普（Carnap）、赖欣巴赫（Reichenbach）和石里克等哲学家，也密切关注他们那个时代的物理学发展。那时，有些人甚至相信双方互动与相互启发的必要性或可能性，或者至少相信和平共存的可能性。如今，这些日子似乎已经过去。

% S30.4.P02
几十年来，物理学家、诺贝尔奖获得者理查德·费曼（Richard Feynmann）的一句话被反复引用：“科学哲学对于科学家的用处，大致就像鸟类学对于鸟的用处一样。”同样获得过诺贝尔物理学奖的史蒂文·温伯格（Steven Weinberg），于1993年写下“哲学不可思议的无效性”（The Unreasonable Ineffectiveness of Philosophy），显然是在俏皮地呼应尤金·维格纳（Eugene Wigner）对“数学不可思议的有效性”的赞叹。在 Weinberg（1993）中，他以一个单独的章节讨论这个问题，并冠以醒目的标题“反对哲学”（Against Philosophy）。天文学家斯蒂芬·霍金（Stephen Hawking）一度宣告哲学已死；较近时期，著名物理学家 Laurence Kraus 也加入了这场争论。2015年，温伯格以《解释世界：现代科学的发现》（To Explain the World: The Discovery of Modern Science）再次引起轩然大波，因而不得不面对历史学家斯蒂芬·谢平（Stephen Shapin）的严厉训诫；后者在发表于《华尔街日报》（The Wall Street Journal）的书评“科学家为什么不该写历史”（Why Scientists Should Not Write History，2015）中提出了批评。这样的清单可以轻而易举地延长，并用大量历史事例加以注释：从歌德（Goethe）对牛顿《光学》（Optics）的否定与拒斥，到20世纪初爱因斯坦与柏格森（Bergson）关于时间本性的激烈论战。不过，即使那些不那么热衷论战的物理学家，通常也接受某种消除式还原论（eliminative reductionism），它消除许多哲学概念和“感受质”（qualia），或赋予它们一种特定的抽象描述，或纯粹数学化的描述。这些概念包括空间、时间、运动、因果性、意向性（intentionality）、目的论（teleology），也包括意义、心灵、自由意志、意识、人格同一性，以及通常所有与形而上学或宗教有关的观念。

% S30.4.P03
幸运的是，统计学与科学哲学之间的关系，并不具有同样的敌意，也没有以耗时的论战为特征。不过，几十年来，两者的对话始终有些困难。这个问题的形成有多个因素，这里只限于其中一个关键方面：自1925年前后起源以来，科学哲学就在发现语境与证成语境之间作出了彻底区分，第30.2节曾简要提及这一点。在逻辑实证主义的启发下，科学哲学最初希望对（统一的）科学进行理性重构，并制定相应的划界标准；因此，它把重点放在证成语境上，而不是研究实践和建设性方法论上。虽然许多科学哲学家几乎立即认识到概率革命的重要性，却悖论性地形成了一种延续至今的分离。正如下节所要勾勒的那样，“追求真理”获得了一种分层结构：由（经验）学科、（研究）方法论和认识论／科学哲学组成的三联画。思想史中的许多争论，以及当代科学与哲学关系中的瓶颈与挑战，都可以用这些“层次”及其相互关系来表达。统计学与哲学目前的关系也同样如此。

% S30.4.P04
这一切都可以在美国贝叶斯统计学家安德鲁·盖尔曼（Andrew Gelman）的工作中看到。在他与人合著的文章“哲学与贝叶斯统计实践”（Philosophy and the Practice of Bayesian Statistics）（Gelman 和 Shalizi，2013）中，两位作者强调，哲学对于统计研究和统计实践都极其重要。他们认为，在这方面，某些统计学家经常表现得“随心所欲”，带有业余作风，而盖尔曼也坦率地把自己列在其中。虽然上述文章把范围限于波普尔、库恩等科学哲学的代表人物，但他们仍鼓励同行与较近时期的哲学研究建立联系。

% S30.4.P05
与此同时，他们又向前迈出一步。Gelman 和 Shalizi（2013）旨在为贝叶斯统计发展一种新哲学，使其在经验上是恰当的，并且能够更好地说明现有统计实践。作者指出，当今统计学惯常把统计方法或假设，与哲学范式或立场联系起来，甚至锚定于其中。他们特别提及归纳主义（inductivism）和假说—演绎模型（hypothetical-deductive model，HD模型）这两种科学观。随后，他们把“经典”统计学的多项成就概括为假说—演绎推理的表现形式：Ronald Fisher 的显著性检验、Egon Pearson 与 Jerzy Neyman 发展的假设检验，以及最后同样由 Neyman 奠基的置信区间理论。

% S30.4.P06
根据 Gelman 和 Shalizi（2013），贝叶斯统计被广泛看作归纳主义在统计学中的一种体现。两种理论被描绘为相互对立、甚至相互冲突的科学观。为了推动一种在经验上恰当的贝叶斯统计新哲学，作者批评了前面勾勒的局面，并主张，贝叶斯统计应被视为一种极其成功的 HD模型；与前述“经典”统计学的表现形式相比，它并不更加归纳，也不更少归纳。此外，他们断言，最成功的贝叶斯推理应用超出了标准看法。在标准看法中，推理被理解为利用以新数据为基础的似然函数，将先验分布更新为后验分布，随后相应调整“信念程度”。实际上，更应把它理解为发展多个模型，即模型拟合，然后评价这些模型，即模型检验，再使其经受一系列严格检验，最终留下最强的模型。因为在作者看来，所有模型（几乎）仍然是“错的”，也就是说，并不包含产生现有数据的概率分布。这里更像是一种试错方法；从历史上说，它与波普尔相联系，而在近期文献中，其最显著的倡导者是 Mayo（1996，2010）。

% S30.4.P07
然而，Gelman 和 Shalizi（2013）没有注意到，归纳主义与 HD模型并非“清楚而分明”（claire et distinct），也很难像他们所暗示的那样，被视为在统计学中发挥作用的两种相互对立的科学观（Starmans，2013）。这种区分更反映了思想家在追求知识的过程中，与推理及其“反转”所作的斗争。归纳尤其经历了漫长的发展过程，具有不同的面貌和侧面，并与相关人物的思想立场紧密相连。这既适用于亚里士多德、培根（Bacon）、休谟、康德（Kant）和密尔（Mill），也适用于波普尔及贝叶斯主义者。这个概念经历了显著转变。严格的主张，例如把朴素归纳主义当作一种逻辑或发现方式，如今已不会再有人提倡；而由许多确证理论支持的温和而“精致”的版本，则只有少数人会认为有问题。HD模型不能声称拥有同样令人尊敬的历史，也不应被看作对归纳的一次简单而成功的修补尝试或替代方案。它过于一般，而且实际上始终与归纳相联系，甚至在弗朗西斯·培根（Francis Bacon）的朴素版本中也是如此。作为一种方法，它并未包含显著的新要素或新洞见，并未解决归纳的问题，而且面临各种历史性与系统性反对意见，其中包括著名的奎因—迪昂论题（Quine-Duhem these）。粗略地说，这一论题主张，不可能为一个单独的、或多或少孤立的假说设计“关键检验”。实际上，假说是其他（辅助）假说、初始条件以及“其他条件不变”（ceteris paribus）条款所组成的网络的一部分，因此，究竟什么受到了检验或被作了归类，并不清楚。

% S30.4.P08
无论如何，把某个人或某种思想归入这两个“阵营”中的任何一方，似乎都免不了受到虚假对立谬误与稻草人谬误的欺骗。此外，包罗一切的综合性认识论理论的时代早已结束，因此，在某种统计方法与某一哲学范式之间寻找或假定一一对应，也就不那么理所当然。尽管存在上述弱点，Gelman 和 Shalizi（2013）试图在统计方法论与认识论之间建立联盟，仍然极具意义和启发性。而且，他们弥合鸿沟的努力并未受到忽视，部分原因是他们并非只面向熟悉这些问题的同行，而是试图触及更广泛的读者；盖尔曼也通过博客与评论，活跃于许多网站。这种努力的重要性十分显著，超出了狭义统计学的范围，而且当前局面绝不能仅归咎于统计学。更应当把它看作某个更一般问题的表现，下一节将使这一点更加清楚。
